% generated by 作成スクリプト/make_steps.py; do not edit
\begin{calc}[手順 1：$\tau\Sigma$ と $\tau\Sigma+\Omega$（$\times10^{-5}$）]
\[ \tau\Sigma = \begin{pmatrix}9.695 & 0.647 & 0.519\\0.647 & 1.082 & 0.239\\0.519 & 0.239 & 1.383\end{pmatrix}\times10^{-5},\qquad \tau\Sigma+\Omega = \begin{pmatrix}19.390 & 0.647 & 0.519\\0.647 & 2.164 & 0.239\\0.519 & 0.239 & 2.765\end{pmatrix}\times10^{-5} \]
例：$(\tau\Sigma)_{11} = 0.025\times0.0038780 = 0.00009695$。$\Omega$ は対角だけなので、対角だけが 2 倍になる：$0.00009695 + 0.00009695 = 0.00019390$。
\end{calc}
\begin{calc}[手順 2：逆行列 $(\tau\Sigma+\Omega)^{-1}$]
\[ (\tau\Sigma+\Omega)^{-1} = \begin{pmatrix}5{,}229.2 & -1{,}470.4 & -854.8\\-1{,}470.4 & 47{,}076.0 & -3{,}789.5\\-854.8 & -3{,}789.5 & 36{,}652.1\end{pmatrix} \]
検算：元の行列と掛けると単位行列になる。たとえば 1 行目 $\times$ 1 列目は
\begin{align*} &19.390\times10^{-5}\times5{,}229.2 + 0.647\times10^{-5}\times(-1{,}470.4) + 0.519\times10^{-5}\times(-854.8)\\ &= 1.0140 + (-0.0095) + (-0.0044) = 1.0000. \end{align*}
\end{calc}
\begin{calc}[手順 3：$M = \tau\Sigma\,(\tau\Sigma+\Omega)^{-1}$]
1 行目 $\times$ 1 列目（$M_{11}$）：
\begin{align*} M_{11} &= 9.695\times10^{-5}\times5{,}229.2 + 0.647\times10^{-5}\times(-1{,}470.4) + 0.519\times10^{-5}\times(-854.8)\\ &= 0.5070 + (-0.0095) + (-0.0044) = 0.493. \end{align*}
9 個の成分を同じように計算すると
\[ M = \begin{pmatrix}0.493 & 0.143 & 0.083\\0.016 & 0.491 & 0.041\\0.012 & 0.052 & 0.493\end{pmatrix} \]
\end{calc}
\begin{calc}[手順 4：$\bm\mu = \bm\pi + M(\bm q-\bm\pi)$]
ビューと均衡の差 $\bm q-\bm\pi$：\textsf{SANA} $0.0039229 - 0.0033040 = 0.0006189$、\textsf{AIV} $0.0005724 - 0.0005722 = 0.0000003$、\textsf{JACK} $0.0006489 - 0.0008376 = -0.0001887$。
\begin{align*}
\mu_1 &= 0.0033040 + (0.493\times0.0006189 + 0.143\times0.0000003 + 0.083\times(-0.0001887))\\
&= 0.0033040 + (0.0003051 + 0.0000000 + (-0.0000156)) = 0.0033040 + 0.0002895 = \mathbf{0.0035935}\\[4pt]
\mu_2 &= 0.0005722 + (0.016\times0.0006189 + 0.491\times0.0000003 + 0.041\times(-0.0001887))\\
&= 0.0005722 + (0.0000098 + 0.0000001 + (-0.0000077)) = 0.0005722 + 0.0000022 = \mathbf{0.0005744}\\[4pt]
\mu_3 &= 0.0008376 + (0.012\times0.0006189 + 0.052\times0.0000003 + 0.493\times(-0.0001887))\\
&= 0.0008376 + (0.0000073 + 0.0000000 + (-0.0000931)) = 0.0008376 + (-0.0000857) = \mathbf{0.0007519}
\end{align*}
\end{calc}
